% !TEX program = xelatex
%
% Bilingual smoke test for this skill. It deliberately exercises every check the
% skill asks for: mixed CJK/Latin text, math, a plot, a table, code, cross
% references, SI units, and a Biber bibliography.
%
% Compile with scripts/cold-build.sh (which cleans first and then verifies the
% log and the embedded fonts), or by hand from this directory:
%
%   latexmk -C main.tex && rm -f main.bbl
%   latexmk -xelatex -synctex=1 -interaction=nonstopmode -halt-on-error main.tex
%
\documentclass[UTF8,fontset=none,11pt]{ctexart}

% --- Font binding: platform decision, not an accident -----------------------
% macOS: pin the two native families XeTeX can open (see assets/macos/).
% Elsewhere: TeX Live's Fandol fonts, which ship with texlive-lang-chinese and
% never touch the desktop font stack.
\usepackage{fontspec}
\IfFontExistsTF{Songti SC Regular}
  {\IfFileExists{../macos/macos-native-fonts.tex}
     {% macOS native CJK fonts for ctex/xeCJK — verified on macOS 26.6.2 (Apple Silicon).
%
% Use with:
%   \documentclass[UTF8,fontset=none]{ctexart}
%   \input{path/to/macos-native-fonts}
%
% Why not [fontset=mac]: CTeX probes STFangsong and calls mktextfm when it is not
% found, and PingFang SC is enumerated by the system font service but XeTeX cannot
% reliably open it by family name. Do not copy system fonts into ~/Library/Fonts to
% silence those warnings — pin the two families that XeTeX can actually open.

\setCJKmainfont{Songti SC Regular}[
  BoldFont      = Songti SC Bold,
  AutoFakeSlant = 0.2
]

\setCJKsansfont{Heiti SC Light}[
  BoldFont   = Heiti SC Medium
]

\setCJKmonofont{Heiti SC Light}[
  BoldFont   = Heiti SC Medium
]

\setCJKfamilyfont{zhsong}{Songti SC Regular}[
  BoldFont   = Songti SC Bold
]

\setCJKfamilyfont{zhhei}{Heiti SC Light}[
  BoldFont   = Heiti SC Medium
]

% No stably accessible separate Kai/FangSong on this macOS build: fall back to
% families that are known to work instead of registering new fonts.
\setCJKfamilyfont{zhkai}{Songti SC Regular}[
  BoldFont   = Songti SC Bold
]

\setCJKfamilyfont{zhfs}{Songti SC Light}[
  BoldFont   = Songti SC Regular
]

\providecommand{\songti}{\CJKfamily{zhsong}}
\providecommand{\heiti}{\CJKfamily{zhhei}}
\providecommand{\kaishu}{\CJKfamily{zhkai}}
\providecommand{\fangsong}{\CJKfamily{zhfs}}
}
     {\setCJKmainfont{Songti SC Regular}}}
  % Fandol ships inside TEXMF: address the files, not the family. On Ubuntu's TeX
  % Live the family name "FandolSong-Regular" does not resolve, the file does.
  {\setCJKmainfont{FandolSong-Regular.otf}[BoldFont={FandolSong-Bold.otf},
                                        ItalicFont={FandolKai-Regular.otf}]
   \setCJKsansfont{FandolHei-Regular.otf}[BoldFont={FandolHei-Bold.otf}]
   \setCJKmonofont{FandolFang-Regular.otf}}

\usepackage[margin=2.5cm]{geometry}
\usepackage{amsmath,amssymb}
\usepackage{booktabs}
\usepackage{pgfplots}
\pgfplotsset{compat=1.18}
\usepackage{siunitx}
\usepackage{listings}
\usepackage{xcolor}
\usepackage{hyperref}
\usepackage{cleveref}
\usepackage[backend=biber,style=numeric]{biblatex}
\addbibresource{refs.bib}

\lstset{
  basicstyle=\ttfamily\small,
  frame=single,
  breaklines=true,
  columns=fullflexible,
}

\title{中英文排版冒烟测试\\[2pt]\large A Bilingual Typesetting Smoke Test}
\author{Toolchain Verification}
\date{\today}

\begin{document}
\maketitle

\section{引言}\label{sec:intro}
中英文混排要求 CJK 字体能够被 XeTeX 直接打开，而不是仅被系统字体服务枚举
\cite{knuth1984,lamport1994}。本节同时验证交叉引用：公式见 \cref{eq:energy}，
图表见 \cref{fig:profile,tab:params}，方法见 \cref{sec:method}。

\section{方法与公式}\label{sec:method}
一维线性对流方程的离散能量变化可写为
\begin{equation}\label{eq:energy}
  \frac{\mathrm{d}}{\mathrm{d}t}\int_{\Omega} \tfrac{1}{2} u^{2}\,\mathrm{d}x
  = -\oint_{\partial\Omega} \tfrac{1}{2} u^{3} n_{x}\,\mathrm{d}s ,
\end{equation}
其中通量项的量级用国际单位表示约为 \qty{1.2}{\metre\per\second}，
无量纲数取 \num{3.14e-2}。

\section{图表与代码}\label{sec:results}
\begin{table}[htbp]
  \centering
  \caption{参数取值 / Parameter values}\label{tab:params}
  \begin{tabular}{lS[table-format=2.2]l}
    \toprule
    参数 & {数值} & 单位 \\
    \midrule
    CFL  & 0.45 & -- \\
    $\Delta x$ & 0.01 & \si{\metre} \\
    \bottomrule
  \end{tabular}
\end{table}

\begin{figure}[htbp]
  \centering
  \begin{tikzpicture}
    \begin{axis}[
        width=0.72\linewidth, height=5cm,
        xlabel={$x$}, ylabel={$u(x)$},
        legend style={font=\small},
        grid=major,
      ]
      \addplot[thick, domain=0:1, samples=80] {sin(deg(pi*x))};
      \addlegendentry{初始条件}
      \addplot[thick, dashed, domain=0:1, samples=80] {sin(deg(pi*(x-0.2)))};
      \addlegendentry{$t=0.2$}
    \end{axis}
  \end{tikzpicture}
  \caption{随时间平移的波形 / Advected profile}\label{fig:profile}
\end{figure}

\begin{lstlisting}[language=Python,caption={最小示例}]
def advect(u, cfl):
    return u - cfl * (u - np.roll(u, 1))
\end{lstlisting}

\section{结论}
若 \cref{tab:params,fig:profile} 正常编号、\cref{eq:energy} 可点击跳转、
参考文献 \cite{knuth1984,lamport1994} 正常解析，且构建日志中没有未定义引用与缺字，
则这套工具链可用于真实论文。

\printbibliography
\end{document}
